\section{Булевы грамматики}
\label{sec:BooleanGrammars}
\tikzsetfigurename{CB_Boolean_}

Дадим определение булевой грамматики.

\begin{definition}
    \textit{Булевой грамматикой} называется $G = (\Sigma,N,P,S)$, где:
    \begin{itemize}
        \item $\Sigma$ и $N$~--- дизъюнктивные конечные непустые множества терминалов и нетерминалов.
        \item $P$~--- конечное множество продукций, каждая вида
        \[
        A\rightarrow \alpha_1\&...\&\alpha_m\&\neg\beta_1\&...\&\neg\beta_n
        \]
        ,где $A \in N, m, n \ge 0, m+n \geq 1$ и $\alpha_i,\beta_j \in (\Sigma \cup N)^*$.
        \item $S \in N$ ~--- стартовый нетерминал.
    \end{itemize}
\end{definition}
\mytodo{Добавить формальное определение языка булевой грамматики: семантика отрицания L(¬β) и отношение выводимости.}
\mytodo{Обсудить проблему немонотонности оператора для булевых грамматик с отрицанием и условие стратификации (stratification), гарантирующее корректность определения.}

Приведем пример булевой грамматики.

\begin{example}
    Следующая булева грамматика порождает язык  $\{a^mb^nc^n\mid m,n \geq 0, m \neq n \}$:
    
    \begin{align*}
    S   &\to A B \& \neg D C \\ 
    A  &\to a A \mid \varepsilon \\ 
    B &\to b B c \mid \varepsilon \\
    C   &\to c C \mid \varepsilon \\ 
    D   &\to aDb \mid \varepsilon
    \end{align*}
    
    Очевидно, что $L(AB) = \{a^mb^nc^n\mid m,n \in \mathbb{N}\}$ и $L(DC) = \{a^nb^nc^m\mid m,n \in \mathbb{N}\}$. Тогда $L(AB)\cap\overline{L(DC)} = \{a^mb^nc^n\mid m,n \geq 0, m \neq n \}$.
\end{example}
\mytodo{Добавить \textbackslash label у примера.}
\mytodo{Добавить пример булевой грамматики с несколькими отрицаниями, демонстрирующий выразительную силу.}

Подробнее о булевых грамматиках можно прочитать в статьях~\cite{Okhotin:2003:BG:1758089.1758123,Okhotin:2014:PMM:2565359.2565379}.

Для матричных алгоритмов синтаксического анализа и поиска путей используется двоичная нормальная форма булевых грамматик~\sidecite{Okhotin:2014:PMM:2565359.2565379,Shemetova2019}.

\begin{definition}[Двоичная нормальная форма булевой грамматики]
    Булева грамматика $G = (\Sigma, N, P, S)$ находится в двоичной нормальной форме, если каждое правило из $P$ имеет один из следующих видов:
    \begin{itemize}
        \item $A \rightarrow B_1 C_1 \& \ldots \& B_m C_m \& \neg D_1 E_1 \& \ldots \& \neg D_k E_k$, где $A,B_i,C_i,D_j,E_j \in N$, $m,k \geqslant 0$ и $m + k \geqslant 1$;
        \item $A \rightarrow a$, где $A \in N$ и $a \in \Sigma$.
    \end{itemize}
\end{definition}
\mytodo{Добавить обсуждение выразительной силы: что добавляет отрицание поверх конъюнкции? Сравнить классы Conj и Bool.}
\mytodo{Добавить обсуждение свойств замкнутости булевых языков.}
\mytodo{Отметить, что лемма о накачке для булевых языков — открытая проблема.}
\mytodo{Добавить обсуждение алгоритмических свойств: сложность синтаксического анализа для булевых грамматик, разрешимость пустоты.}
\mytodo{Добавить примеры практического применения булевых грамматик.}
\mytodo{Сформулировать теорему о существовании двоичной нормальной формы для любой булевой грамматики и добавить пример приведения.}
\mytodo{Добавить сравнение конъюнктивных и булевых грамматик с MCFG: ортогональные расширения КС-грамматик, соотношение выразительной силы.}
\mytodo{Добавить иллюстрации: диаграммы вывода, соотношение классов языков.}
